\section{小耦合下威尔逊流的性质}

本节的目的部分在于表明威尔逊流可以直接在微扰论中研究，部分在于核验：在正流时间
计算的局域规范不变观测量的期望值确为重正化的量。

为简单起见，微扰展开在连续理论中借助维数正规化进行讨论。取规范群为 $\SUn$，并
假设有 $\Nf$ 味无质量夸克。作为代表性例子，考虑可观测量
\begin{equation}
   E=\frac{1}{4}G_{\mu\nu}^aG_{\mu\nu}^a
   \label{eq:Edens}
\end{equation}
并将其期望值计算到规范耦合的次领头阶。


\subsection{规范固定}

流方程~\eqref{eq:floweq} 在不依赖 $t$ 的规范变换下不变。因此 $E$ 是基本场
$A_{\mu}$ 的规范不变函数，其期望值从而可在任意规范中计算。

在微扰论中，流方程的规范不变性带来一些技术上的复杂之处；通过考虑修改后的方程
\begin{equation}
  \dot{B}_{\mu}=D_{\nu}G_{\nu\mu}+\lambda D_{\mu}\partial_{\nu}B_{\nu}.
  \label{eq:modflow}
\end{equation}
更容易避开它们。对规范参数 $\lambda$ 的任意给定值，方程~\eqref{eq:modflow} 的解与
$\lambda=0$ 处的解通过
\begin{equation}
  B_{\mu}=\Lambda\!\left.B_{\mu}\right|_{\lambda=0}\Lambda^{-1}
  +\Lambda\partial_{\mu}\Lambda^{-1},
  \label{eq:gaugetrafo}
\end{equation}
相联系，其中规范变换 $\Lambda(t,x)$ 由
\begin{equation}
  \dot{\Lambda}=-\lambda \partial_{\nu}B_{\nu}\Lambda,
  \qquad\left.\Lambda\right|_{t=0}=1.
  \label{eq:Lambdatrafo}
\end{equation}
确定。于是可以用修改后的流方程计算 $E$ 的期望值。而且可以自由地取 $\lambda=1$，
事实证明这是一个特别方便的选择。

注意，使用修改后的流方程并不妨碍基本场的规范固定，因为 $E$ 不变，因而仍是后者的
规范不变函数。

\subsection{修改后流方程的解}

在微扰论中，将规范势按裸耦合重新标度
\begin{equation}
  A_{\mu}\to g_0A_{\mu},
\end{equation}
然后把泛函积分按 $g_0$ 的幂次展开。于是流 $B_{\mu}(t,x)$ 成为耦合的函数，具有如下
形式的渐近展开
\begin{equation}
  B_{\mu}=\sum_{k=1}^{\infty}g_0^kB_{\mu,k},
  \qquad
  \left.B_{\mu,k}\right|_{t=0}=\delta_{k1}A_{\mu}.
  \label{eq:Bseries}
\end{equation}
把该级数代入方程~\eqref{eq:modflow}，并取 $\lambda=1$，便得到一列方程
\begin{equation}
  \dot{B}_{\mu,k}-\partial_{\nu}\partial_{\nu}B_{\mu,k}=
  R_{\mu,k},
  \quad k=1,2,\ldots,
  \label{eq:tower}
\end{equation}
其右端表达式为
\begin{align}
  R_{\mu,1}&=0,
  \\
  R_{\mu,2}&=2[B_{\nu,1},\partial_{\nu}B_{\mu,1}]
            -[B_{\nu,1},\partial_{\mu}B_{\nu,1}],
  \\
  R_{\mu,3}&=2[B_{\nu,2},\partial_{\nu}B_{\mu,1}]
            +2[B_{\nu,1},\partial_{\nu}B_{\mu,2}]
            \nonumber\\
  &\hphantom{{}=}
            -[B_{\nu,2},\partial_{\mu}B_{\nu,1}]
            -[B_{\nu,1},\partial_{\mu}B_{\nu,2}]
            +[B_{\nu,1},[B_{\nu,1},B_{\mu,1}]],
\end{align}
如此等等。特别地，在 $D$ 维情形，领头阶方程给出
\begin{align}
  B_{\mu,1}(t,x)&=\int\rmd^Dy\,K_t(x-y)A_{\mu}(y),
  \label{eq:B1ker}\\
  K_t(z)&=\int{\rmd^Dp\over(2\pi)^D}\,\rme^{ipz}\rme^{-tp^2}=
  {\rme^{-z^2/4t}\over(4\pi t)^{D/2}},
  \label{eq:Kt}
\end{align}
这明确表出流是一种光滑化操作。更确切地说，规范势是在空间的一个球形范围内被平均
的，该范围在四维时的均方半径等于 $\sqrt{8t}$。

高阶方程~\eqref{eq:tower} 可以逐一求解，只需注意
\begin{equation}
   B_{\mu,k}(t,x)=\int_0^t\rmd s\int\rmd^Dy\,K_{t-s}(x-y)R_{\mu,k}(s,y).
   \label{eq:Bksol}
\end{equation}
回顾式~\eqref{eq:B1ker},\,\eqref{eq:Kt} 即清楚，此公式生成树状表达式，基本场
$A_{\mu}$ 附着在树的端点上。


\subsection{$\langle E\rangle$ 的展开}

把级数~\eqref{eq:Bseries} 代入
\begin{equation}
\begin{split}
  \langle E\rangle={}&
           \frac{1}{2}\langle\partial_{\mu}B^a_{\nu}\partial_{\mu}B^a_{\nu}-
           \partial_{\mu}B^a_{\nu}\partial_{\nu}B^a_{\mu}\rangle\\
  &+f^{abc}\langle\partial_{\mu}B^a_{\nu}B^b_{\mu}B^c_{\nu}\rangle
           +\frac{1}{4}f^{abe}f^{cde}\langle
           B^a_{\mu}B^b_{\nu}B^c_{\mu}B^d_{\nu}\rangle,
\end{split}
\label{eq:Eexpansion}
\end{equation}
便得到一串阶数递增的项。最低阶项是
\begin{equation}
  {\cal E}_0=\frac{1}{2}g_0^2
             \langle\partial_{\mu}B^a_{\nu,1}\partial_{\mu}B^a_{\nu,1}-
             \partial_{\mu}B^a_{\nu,1}\partial_{\nu}B^a_{\mu,1}\rangle
  \label{eq:E0def}
\end{equation}
而次一阶的项是
\begin{align}
  {\cal E}_1&=g_0^3
    f^{abc}\langle\partial_{\mu}B^a_{\nu,1}B^b_{\mu,1}B^c_{\nu,1}\rangle,
  \\
  {\cal E}_2&=
    g_0^3
    \langle\partial_{\mu}B^a_{\nu,2}\partial_{\mu}B^a_{\nu,1}-
           \partial_{\mu}B^a_{\nu,2}\partial_{\nu}B^a_{\mu,1}\rangle.
  \label{eq:E2def}
\end{align}
每一项都是规范耦合的幂级数；只要把系数 $B_{\mu,k}(t,x)$ 用基本场 $A_{\mu}(x)$ 表出，
并按标准费曼规则展开后者的关联函数，便可将其求出。

实际计算中，通过插入傅里叶表示
\begin{align}
  B^a_{\mu,1}(t,x)&=\int_p\rme^{ipx}\rme^{-tp^2}\tilde{A}^a_{\mu}(p),
  \label{eq:B1fourier}\\
  B^a_{\mu,2}(t,x)&=if^{abc}\int_0^t\rmd s\int_{q,r}
  \rme^{i(q+r)x}\rme^{-s(q^2+r^2)-(t-s)(q+r)^2}
  \nonumber\\
  &\hphantom{{}=}\times
  \bigl\{\delta_{\mu\lambda}r_{\sigma}-\delta_{\mu\sigma}q_{\lambda}
  +\frac{1}{2}\delta_{\sigma\lambda}(q-r)_{\mu}\bigr\}
  \tilde{A}_{\sigma}^b(q)\tilde{A}_{\lambda}^c(r),
  \label{eq:B2fourier}
\end{align}
以及高阶场 $B_{\mu,k}(t,x)$ 的相应表达式而过渡到动量空间是有利的。为简化记号，这里
引入了简写
\begin{equation}
  \int_{p}=\int{\rmd^D p\over(2\pi)^D}
  \label{eq:intp}
\end{equation}

注意，式~\eqref{eq:B2fourier} 中的流时间积分类似于费曼参数积分。特别地，若对胶子
关联函数各图的贡献使用费曼参数，最终得到的高斯动量积分在任何维数 $D$ 都容易求值。
对这些参数的积分与通常遇到的几乎一样，唯一区别是流时间参数只积到 $t$ 而非无穷。


\subsection{${\cal E}_0$ 的计算}

对最低阶项~\eqref{eq:E0def}，上一小节概述的步骤给出公式
\begin{equation}
  {\cal E}_0=\frac{1}{2}g_0^2(N^2-1)
  \int_p\rme^{-2tp^2}
  (p^2\delta_{\mu\nu}-p_{\mu}p_{\nu})
  D(p)_{\mu\nu},
  \label{eq:E0int}
\end{equation}
其中 $D(p)_{\mu\nu}$ 表示未重正化的完整胶子传播子。取 $D=4-2\eps$
并选费曼规范，传播子具有形式
\begin{align}
  D(p)_{\mu\nu}&={1\over(p^2)^2}\left\{
  (p^2\delta_{\mu\nu}-p_{\mu}p_{\nu})(1-\omega(p))^{-1}+p_{\mu}p_{\nu}
  \right\},
  \\
  \omega(p)&=\sum_{k=1}^{\infty}g_0^{2k}(p^2)^{-k\eps}\omega_k,
\end{align}
由此推得
\begin{equation}
  {\cal E}_0=\frac{1}{2}g_0^2{N^2-1\over(8\pi t)^{D/2}}(D-1)
  \left\{1+g_0^2(2t)^{\eps}{\Gamma(2-2\eps)\over\Gamma(2-\eps)}\omega_1
  +\ldots\right\}.
  \label{eq:E0res}
\end{equation}
这一步的计算很容易完成，只需引用胶子自能的已知结果
\begin{equation}
  \omega_1={1\over 16\pi^2}(4\pi\rme^{-\euler})^{\eps}
  \left\{N\left(\frac{5}{3\eps}+\frac{31}{9}\right)
  -\Nf\left(\frac{2}{3\eps}+\frac{10}{9}\right)
  +\rmO(\eps)\right\}
  \label{eq:omega1}
\end{equation}
其中 $\euler=0.577\ldots$ 为欧拉常数。


\subsection{$\langle E\rangle$ 到 $g_0^4$ 阶的计算}

在次领头阶，$E$ 的期望值接收来自 ${\cal E}_0$ 以及由流方程展开生成的、直到
$g_0^4$ 阶的项 ${\cal E}_1,{\cal E}_2,\ldots$ 的贡献（参见 2.3 小节）。后者中有
一些涉及三点胶子顶点，但在所有情形下都只需要领头阶的胶子关联函数。

以项 ${\cal E}_2$ 为例，计算的纯代数部分给出积分
\begin{equation}
   {\cal E}_2=N(N^2-1)g_0^4\int_0^t\rmd s\int_{q,r}
   {\rme^{-(2t-s)p^2-s(q^2+r^2)}\over p^2q^2r^2}
   \nonumber\\
   \quad\times\bigl\{
   (D-1)p^2(p^2+q^2+r^2)+2(D-2)(q^2r^2-(qr)^2)\bigr\}
   +\rmO(g_0^6)
   \label{eq:E2int}
\end{equation}
其中 $p=q+r$。这个积分看似相当复杂，但被积函数分子中的多项式使表达式可以化简并
最终解析求值（进一步的细节见附录~\ref{app:B}）。计算结果
\begin{equation}
   {\cal E}_2=N(N^2-1){g_0^4\over(4\pi)^D(2t)^{D-2}}
   \bigl\{\frac{9}{2\eps}-\frac{3}{2}+45\ln2-\frac{45}{2}\ln3+
   \rmO(\eps)\bigr\}+\rmO(g_0^6),
   \label{eq:E2res}
\end{equation}
表明这些贡献并非没有紫外奇异性。

其余项的计算遵循同样的模式，不会出现任何额外的困难。汇集全部贡献即得结果
\begin{align}
  \langle E\rangle&=\frac{1}{2}g_0^2{N^2-1\over(8\pi t)^{D/2}}(D-1)
  \bigl\{1+c_1g_0^2+\rmO(g_0^4)
  \bigr\},
  \label{eq:Epert}\\
  c_1&={1\over16\pi^2}(4\pi)^{\eps}(8t)^{\eps}\bigl\{
  N\left(\frac{11}{3\eps}+\frac{52}{9}-3\ln3\right)
  -\Nf\left(\frac{2}{3\eps}+\frac{4}{9}-\frac{4}{3}\ln2\right)
  +\rmO(\eps)\bigr\}.
\end{align}


\subsection{重正化}

裸耦合 $g_0$ 与 $\MSbar$ 方案中的重正化耦合 $g$ 及相应的归一化质量 $\mu$
按文献~\cite{BardeenEtAl} 相联系：
\begin{align}
  g_0^2&=g^2\mu^{2\eps}(4\pi\rme^{-\euler})^{-\eps}
  \Bigl\{1-\frac{1}{\eps}b_0g^2+\rmO(g^4)\Bigr\},
  \label{eq:g0rel}\\
  b_0&={1\over16\pi^2}\left\{\frac{11}{3}N-\frac{2}{3}\Nf\right\}.
\end{align}
现在若把式~\eqref{eq:g0rel} 写成重正化耦合的展开，则与 $1/\eps$ 成比例的项相消，
在 $\eps=0$ 处得到
\begin{align}
  \langle E\rangle&={3(N^2-1)g^2\over128\pi^2t^2}
  \bigl\{1+\bar{c}_1g^2+\rmO(g^4)
  \bigr\},
  \label{eq:Eren}\\
  \bar{c}_1&=
  {1\over16\pi^2}\bigl\{
  N\left(\frac{11}{3}L+\frac{52}{9}-3\ln3\right)
  -\Nf\left(\frac{2}{3}L+\frac{4}{9}-\frac{4}{3}\ln2\right)\bigr\},
\end{align}
其中 $L=\ln(8\mu^2t)+\euler$。在这一阶规范耦合下，$\langle E\rangle$
结果证明是一个无需重正化的量，因此编码了理论的某种物理性质。

用标度 $q=(8t)^{-1/2}$ 处的跑动耦合 $\alpha(q)$ 表出时，展开式~\eqref{eq:Eren}
具有形式
\begin{align}
  \langle E\rangle&={3(N^2-1)\over32\pi t^2}\alpha(q)
  \bigl\{1+k_1\alpha(q)+\rmO(\alpha^2)
  \bigr\},
  \label{eq:Erun}\\
  k_1&=
  {1\over4\pi}\bigl\{
  N\left(\frac{11}{3}\euler+\frac{52}{9}-3\ln3\right)
  -\Nf\left(\frac{2}{3}\euler+\frac{4}{9}-\frac{4}{3}\ln2\right)\bigr\}.
\end{align}
特别地，对 $N=3$ 得到
\begin{equation}
  \langle E\rangle={3\over4\pi t^2}\alpha(q)
  \bigl\{1+k_1\alpha(q)+\rmO(\alpha^2)
  \bigr\},
  \qquad
  k_1=1.0978+0.0075\times\Nf.
  \label{eq:E3}
\end{equation}
此情形下次领头阶修正相当小；当 $\Nf\leq3$ 时，它大致相当于动量 $q$
改变一个约 $2$ 的因子。



